% =============================================================
%  CS102-F26 Digital Design and Computer Architecture I
%  Lecture 14 -- Timing of Sequential Circuits (Unit 14)
%  Follows Section 5.5 (5.5.1 and 5.5.2) of
%  A Student's Computer Architecture Book I.
%
%  Build the slides and the print handout:  npm run materials -- lecture-14
% =============================================================
\makeatletter
\def\input@path{{templates/latex/}{../}}
\makeatother

\documentclass[10pt]{wosm-slides}
\usepackage{cs102-lectures}[2026/10/20]

\title[Lecture 14 \textperiodcentered\ Timing of Sequential Circuits]{Timing of
  Sequential Circuits}
\subtitle{Lecture 14: setup time, hold time and how fast the clock can go}

% ---------- Helpers for this lecture --------------------------
\tikzset{
  sig/.style={draw=wosm-text, line width=0.9pt},
  dim/.style={draw=wosm-text, line width=0.5pt,
    {Stealth[length=3pt]}-{Stealth[length=3pt]}},
  span/.style={draw=wosm-text, line width=0.6pt, fill=wosm-surface,
    minimum height=0.42cm, inner sep=0pt, font=\scriptsize},
  logic/.style={draw=wosm-text, line width=0.6pt, minimum width=2cm,
    minimum height=1.1cm, font=\scriptsize, align=center},
}
% \spanbox{<from>}{<to>}{<y>}{<label>}: one stretch of time on a timeline.
\newcommand*\spanbox[4]{%
  \node[span, minimum width={(#2-#1)*1.1cm}, anchor=west] at (#1,#3) {#4};}

% Two flip-flops with logic in between, as in Figures 5.14 and 5.15.
% \twoflops{<label inside the logic block>}
\newcommand*\twoflops[1]{%
  \begin{tikzpicture}[gates]
    \pic (s) at (0,0) {dff};
    \pic (t) at (5,0) {dff};
    \node[logic] (c) at (2.5,0.27) {#1};
    \node[above, font=\scriptsize] at (0,0.55) {source};
    \node[above, font=\scriptsize] at (5,0.55) {target};
    \draw[wire] (s-d) -- ++(-0.5,0) node[left] {$d$};
    \draw[wire] (s-q) -- (c.west);
    \draw[wire] (c.east) -- (t-d);
    \draw[wire] (t-q) -- ++(0.5,0) node[right] {$q$};
    \draw[wire] (-1.2,-1.05) node[left] {clk} -| ($(t-clk)+(-0.3,0)$) -- (t-clk);
    \draw[wire] ($(s-clk)+(-0.3,-0.78)$) |- (s-clk);
    \fill[wosm-text] ($(s-clk)+(-0.3,-0.78)$) circle (1.5pt);
    \node[font=\scriptsize, text=wosm-brick, anchor=north] at ($(s-clk)+(-0.3,-0.82)$)
      {launch};
    \node[font=\scriptsize, text=wosm-brick, anchor=north] at ($(t-clk)+(-0.3,-0.82)$)
      {capture};
  \end{tikzpicture}}

% =============================================================
\begin{document}

\maketitle

\begin{frame}{Today}
  \tableofcontents
\end{frame}

% =============================================================
\section{Trouble at the edge}

\begin{frame}{What if $D$ changes at the edge?}
  \begin{itemize}
    \item In the master--slave flip-flop \unitref{12}, the rising edge closes the
      master and opens the slave. But what if $D$ changes at that very moment?
  \end{itemize}

  \begin{columns}[T, totalwidth=\textwidth]
    \column{0.48\textwidth}
      \wosmlabel{$D$ changes just after the edge}\par\smallskip
      The master is enabled by $\mathrm{clk}'$, which comes through a NOT gate. For
      a moment after the edge, $\mathrm{clk}'$ is still 1 while clk is already 1:
      \alert{both} latches are transparent, and the new $D$ runs straight through
      to $Q$.
    \column{0.48\textwidth}
      \wosmlabel{$D$ changes just before the edge}\par\smallskip
      The master closes while its gates are still settling on the new $D$. $Q_m$
      can end up 0 or 1, or swing between them for a while: the flip-flop is
      \alert{metastable}, as the SR latch was in Lecture 11.
  \end{columns}

  \medskip
  \begin{itemize}
    \item The book's simulations (Figures 5.11 and 5.12) show both: $Q$ takes the
      wrong value, or oscillates.
    \item The cure: keep $D$ steady for a short while around every rising edge.
  \end{itemize}
\end{frame}

\begin{frame}{Setup time and hold time}
  \begin{center}
    \begin{tikzpicture}[x=0.85cm]
      \fill[wosm-accent!15] (3,-1.12) rectangle (5,-0.53);
      \draw[wosm-rule, line width=0.6pt] (4,0.5) -- (4,-2.4);
      \node[left, font=\scriptsize] at (-0.1,0.17) {clk};
      \draw[sig] (0,0) -- (4,0) -- (4,0.35) -- (9,0.35);
      \node[left, font=\scriptsize] at (-0.1,-0.83) {$D$};
      \draw[sig] (0,-0.65) -- (1.4,-0.65) -- (1.6,-1) -- (6.4,-1) -- (6.6,-0.65)
        -- (9,-0.65);
      \draw[sig] (0,-1) -- (1.4,-1) -- (1.6,-0.65) -- (6.4,-0.65) -- (6.6,-1)
        -- (9,-1);
      \node[font=\scriptsize] at (2.3,-0.83) {$d$};
      \node[left, font=\scriptsize] at (-0.1,-2.03) {$Q$};
      \draw[sig] (0,-1.85) -- (4.7,-1.85) -- (4.9,-2.2) -- (9,-2.2);
      \draw[sig] (0,-2.2) -- (4.7,-2.2) -- (4.9,-1.85) -- (9,-1.85);
      \node[font=\scriptsize] at (6.9,-2.03) {$d$};
      \draw[dim] (3,-1.3) -- node[below, font=\scriptsize] {$t_{\mathrm{setup}}$}
        (4,-1.3);
      \draw[dim] (4,-1.3) -- node[below, font=\scriptsize] {$t_{\mathrm{hold}}$}
        (5,-1.3);
      \draw[dim] (4,-2.45) -- node[below, font=\scriptsize] {$t_{\mathrm{clk2q}}$}
        (4.8,-2.45);
    \end{tikzpicture}
  \end{center}

  \begin{definition}[Timing of a flip-flop]
    \emph{Setup time} $t_{\mathrm{setup}}$: how long $D$ must be stable \emph{before}
    the clock edge. \emph{Hold time} $t_{\mathrm{hold}}$: how long it must stay
    stable \emph{after} it. \emph{Clock-to-Q delay} $t_{\mathrm{clk2q}}$: how long
    after the edge the new value appears at $Q$.
  \end{definition}
\end{frame}

% =============================================================
\section{Setup time constraint}

\begin{frame}{From one flip-flop to the next}
  \begin{center}
    \twoflops{combinational\\ logic\\ $t_{\mathrm{pd}}$, $t_{\mathrm{cd}}$}
  \end{center}

  \begin{itemize}
    \item The basic pattern of every synchronous circuit: a value leaves the
      \emph{source} flip-flop at one clock edge, passes through combinational logic,
      and is stored by the \emph{target} flip-flop at the next edge.
    \item $t_{\mathrm{pd}}$ and $t_{\mathrm{cd}}$ are the logic's propagation and
      contamination delays \unitref{9}: its longest and its shortest path.
    \item The edge reaches the source at time $T_{\mathrm{launch}}$ and the target at
      $T_{\mathrm{capture}}$.
  \end{itemize}
\end{frame}

\begin{frame}{The setup time constraint}
  \begin{center}
    \begin{tikzpicture}[x=1.1cm]
      \foreach \t in {1,7} {\draw[wosm-rule, line width=0.6pt] (\t,0.4) -- (\t,-1);}
      \timingwave{0}{clk}{0,1,1,1,0,0,0,1}
      \spanbox{1}{1.8}{-0.75}{$t_{\mathrm{clk2q}}$}
      \spanbox{1.8}{5.2}{-0.75}{$t_{\mathrm{pd}}$}
      \spanbox{6}{7}{-0.75}{$t_{\mathrm{setup}}$}
      \node[font=\scriptsize, text=wosm-muted] at (5.6,-0.75) {slack};
      \node[below, font=\scriptsize, text=wosm-brick] at (1,-1) {$T_{\mathrm{launch}}$};
      \node[below, font=\scriptsize, text=wosm-brick] at (7,-1)
        {$T_{\mathrm{capture}} + t_{\mathrm{clk}}$};
      \draw[dim] (1,0.55) -- node[above, font=\scriptsize] {$t_{\mathrm{clk}}$} (7,0.55);
    \end{tikzpicture}
  \end{center}

  The new value must be at the target's $D$ at least $t_{\mathrm{setup}}$ before the
  \emph{next} edge, even along the longest path:
  \[
    T_{\mathrm{launch}} + t_{\mathrm{clk2q}} + t_{\mathrm{pd}} + t_{\mathrm{setup}}
    \;\le\; T_{\mathrm{capture}} + t_{\mathrm{clk}}.
  \]
  When the clock reaches both flip-flops at the same time
  ($T_{\mathrm{launch}} = T_{\mathrm{capture}}$):
  \[
    t_{\mathrm{clk}} \;\ge\; t_{\mathrm{clk2q}} + t_{\mathrm{pd}} + t_{\mathrm{setup}}.
  \]
\end{frame}

\begin{frame}{How fast can the clock go?}
  \begin{columns}[T, totalwidth=\textwidth]
    \column{0.48\textwidth}
      \begin{example}
        $t_{\mathrm{clk2q}} = 30$\,ps, $t_{\mathrm{pd}} = 420$\,ps,
        $t_{\mathrm{setup}} = 50$\,ps.\par\smallskip
        $t_{\mathrm{clk}} \ge 30 + 420 + 50 = 500$\,ps, so the clock can run at
        up to $1 / 500\,\mathrm{ps} = 2$\,GHz.
      \end{example}
    \column{0.48\textwidth}
      \begin{itemize}
        \item The slowest path between any two flip-flops sets the clock period of
          the whole circuit.
        \item Too slow? Run the clock slower, make the logic on that path faster,
          or put another register in the middle of it: the idea behind pipelining
          \unitref{20}.
      \end{itemize}
  \end{columns}

  \begin{remark}
    The synchronous counter of Lecture 13: the longest path runs from $y_0$
    through the chain of AND gates to the last $T$ input. The more bits, the
    longer the chain, the slower the clock.
  \end{remark}
\end{frame}

% =============================================================
\section{Hold time constraint}

\begin{frame}{The hold time constraint}
  \begin{center}
    \begin{tikzpicture}[x=1.1cm]
      \draw[wosm-rule, line width=0.6pt] (1,0.4) -- (1,-1.55);
      \timingwave{0}{clk}{0,1,1,1,0,0,0}
      \node[left, font=\scriptsize] at (-0.1,-0.75) {source $\to$ target};
      \spanbox{1}{1.8}{-0.75}{$t_{\mathrm{clk2q}}$}
      \spanbox{1.8}{2.7}{-0.75}{$t_{\mathrm{cd}}$}
      \node[left, font=\scriptsize] at (-0.1,-1.3) {target};
      \spanbox{1}{2.3}{-1.3}{$t_{\mathrm{hold}}$}
      \draw[wosm-brick, dashed, line width=0.5pt] (2.7,-0.5) -- (2.7,-1.55);
      \node[right, font=\scriptsize, text=wosm-brick] at (2.75,-1.3)
        {new value arrives};
      \node[below, font=\scriptsize, text=wosm-brick] at (1,-1.55)
        {$T_{\mathrm{launch}} = T_{\mathrm{capture}}$};
    \end{tikzpicture}
  \end{center}

  The same edge that the target is still capturing also launches the next value from
  the source. Along the \emph{shortest} path, it must not arrive before the target's
  hold time is over:
  \[
    T_{\mathrm{launch}} + t_{\mathrm{clk2q}} + t_{\mathrm{cd}}
    \;\ge\; T_{\mathrm{capture}} + t_{\mathrm{hold}},
    \qquad\text{or, with no skew,}\qquad
    t_{\mathrm{clk2q}} + t_{\mathrm{cd}} \;\ge\; t_{\mathrm{hold}}.
  \]
\end{frame}

\begin{frame}{Fixing a hold violation}
  \begin{itemize}
    \item $t_{\mathrm{clk}}$ does not appear in the hold constraint. A circuit that
      violates it fails at \alert{every} clock speed: slowing the clock does not
      help.
    \item The fix is to make the shortest path slower: add buffers (or other gates
      that keep the value) on it.
    \item Two flip-flops wired directly, with no logic between them, have
      $t_{\mathrm{cd}} = 0$: they need $t_{\mathrm{clk2q}} \ge t_{\mathrm{hold}}$.
  \end{itemize}

  \begin{example}
    $t_{\mathrm{clk2q}} = 30$\,ps, $t_{\mathrm{cd}} = 20$\,ps,
    $t_{\mathrm{hold}} = 60$\,ps. Then $30 + 20 = 50 < 60$: the hold time is
    violated by 10\,ps. Adding at least 10\,ps of delay to the shortest path fixes
    it, as long as the longest path still meets the setup constraint.
  \end{example}
\end{frame}

% =============================================================
\section{Clock skew}

\begin{frame}{When the clock arrives late}
  \begin{itemize}
    \item The clock is a signal on a wire too: it can reach the target later than
      the source. The difference is the \emph{clock skew},
      $t_{\mathrm{skew}} = T_{\mathrm{capture}} - T_{\mathrm{launch}}$.
    \item Put it into the two constraints of the book:
  \end{itemize}
  \begin{align*}
    \text{setup:}\quad & t_{\mathrm{clk}} \;\ge\; t_{\mathrm{clk2q}} + t_{\mathrm{pd}}
      + t_{\mathrm{setup}} - t_{\mathrm{skew}} \\
    \text{hold:}\quad & t_{\mathrm{clk2q}} + t_{\mathrm{cd}} \;\ge\;
      t_{\mathrm{hold}} + t_{\mathrm{skew}}
  \end{align*}
  \begin{itemize}
    \item A late capture edge gives the longest path more time, but the shortest
      path must now wait even longer. Skew helps one constraint and hurts the other.
  \end{itemize}

  \begin{remark}
    The book builds the DX8 without doing these checks, and its epilogue leaves them
    to the reader. You now know how.
  \end{remark}
\end{frame}

\begin{frame}{Try it}
  \begin{enumerate}
    \item $t_{\mathrm{clk2q}} = 40$\,ps, $t_{\mathrm{pd}} = 160$\,ps,
      $t_{\mathrm{setup}} = 50$\,ps. What is the fastest clock?\par
      \uncover<2->{\small $t_{\mathrm{clk}} \ge 250$\,ps: 4\,GHz.}
    \item The same circuit has $t_{\mathrm{cd}} = 15$\,ps and
      $t_{\mathrm{hold}} = 70$\,ps. Does it work at 1\,GHz?\par
      \uncover<3->{\small No: $40 + 15 = 55 < 70$. It fails at any clock speed;
        add at least 15\,ps to the shortest path.}
    \item Now the clock reaches the target 20\,ps after the source. Redo both
      checks.\par
      \uncover<4->{\small Setup: $t_{\mathrm{clk}} \ge 230$\,ps (about 4.35\,GHz).
        Hold: $55 \ge 70 + 20$ fails by 35\,ps.}
    \item Two flip-flops wired directly, $t_{\mathrm{clk2q}} = 30$\,ps,
      $t_{\mathrm{hold}} = 20$\,ps. Is that safe? And with 15\,ps of skew?\par
      \uncover<5->{\small Without skew, yes: $30 \ge 20$. With it, no:
        $30 < 20 + 15$.}
  \end{enumerate}
\end{frame}

% =============================================================
\section{Summary}

\begin{frame}{Summary}
  \begin{itemize}
    \item A flip-flop needs $D$ stable from $t_{\mathrm{setup}}$ before the clock
      edge until $t_{\mathrm{hold}}$ after it; otherwise it can store the wrong value
      or go metastable. Its new value appears $t_{\mathrm{clk2q}}$ after the edge.
    \item Setup: the longest path must fit in one period,
      $t_{\mathrm{clk}} \ge t_{\mathrm{clk2q}} + t_{\mathrm{pd}} + t_{\mathrm{setup}}$.
      This sets the fastest clock.
    \item Hold: the shortest path must not be too quick,
      $t_{\mathrm{clk2q}} + t_{\mathrm{cd}} \ge t_{\mathrm{hold}}$, at any clock
      speed.
    \item Clock skew loosens the setup constraint and tightens the hold constraint.
  \end{itemize}
\end{frame}

\begin{frame}{Before the next lecture}
  \begin{itemize}
    \item Read Section~5.5 of the book.
    \item For the 3-bit synchronous counter of Lecture 13, list every path from a
      flip-flop output to a flip-flop input. Which one is the longest, and which the
      shortest?
    \item Read the beginning of Chapter~6 and Section~6.1.
  \end{itemize}

  \begin{block}{Next lecture \unitref{15}}
    Instruction set architecture: the Von Neumann architecture, and the DX8's
    instruction format and instruction set.
  \end{block}
\end{frame}

\end{document}
